\documentclass[a4paper]{article}
% Español (México) con XeLaTeX: fontspec para fuentes Unicode, polyglossia
% para la división silábica y los nombres del documento en la variante mexicana.
\usepackage{fontspec}
\usepackage{polyglossia}
\setdefaultlanguage[variant=mexican]{spanish}
% Los nombres chinos van como «pinyin (original)»: xeCJK da a los caracteres
% Han su propia fuente; sin ella XeLaTeX los omite con un aviso.
% FandolSong no cubre algunos caracteres de nombres (U+73FA); WenQuanYi Zen Hei sí.
\usepackage[AutoFallBack=true]{xeCJK}
\setCJKmainfont{FandolSong-Regular.otf}
\setCJKfallbackfamilyfont{\CJKrmdefault}{WenQuanYi Zen Hei}
% polyglossia no traduce los nombres de listings.
\AtBeginDocument{\providecommand{\lstlistingname}{}\renewcommand{\lstlistingname}{Listado}%
  \providecommand{\lstlistlistingname}{}\renewcommand{\lstlistlistingname}{Índice de listados}}
\usepackage{amsmath,amssymb}
\usepackage{graphicx}
\usepackage[margin=2.35cm]{geometry}
\usepackage[most]{tcolorbox}
\usepackage{etoolbox}
\usepackage{listings}
\usepackage{xcolor}
\usepackage{hyperref}
\usepackage{booktabs,longtable,tabularx,array}
\usepackage{subcaption}
\usepackage{float}
\usepackage{enumitem}
\usepackage{microtype}
\usepackage{fancyhdr}
\usepackage{needspace}

\definecolor{kbBack}{HTML}{F2F6FA}
\definecolor{kbFrame}{HTML}{9DB4CC}
\definecolor{kbTitle}{HTML}{52708F}
\definecolor{ibBack}{HTML}{FBF5E7}
\definecolor{ibFrame}{HTML}{C9A961}
\definecolor{ibTitle}{HTML}{7D5F1E}
\definecolor{wbBack}{HTML}{FCF2F0}
\definecolor{wbFrame}{HTML}{D6A096}
\definecolor{wbTitle}{HTML}{9E4F43}
\definecolor{dbBack}{HTML}{F4F7F3}
\definecolor{dbFrame}{HTML}{AEC2AC}
\definecolor{dbTitle}{HTML}{5F7F64}

\tcbset{softbox/.style={
  enhanced,breakable,boxrule=0.8pt,arc=2pt,left=3mm,right=3mm,
  top=2mm,bottom=2mm,coltitle=white,fonttitle=\bfseries,
  toptitle=0.6mm,bottomtitle=0.6mm
}}
\newtcolorbox{knowledgebox}[1]{softbox,colback=kbBack,colframe=kbFrame,colbacktitle=kbTitle,title={#1}}
\newtcolorbox{importantbox}[1]{softbox,colback=ibBack,colframe=ibFrame,colbacktitle=ibTitle,title={#1}}
\newtcolorbox{warningbox}[1]{softbox,colback=wbBack,colframe=wbFrame,colbacktitle=wbTitle,title={#1}}
\newtcolorbox{dialoguebox}[1]{softbox,colback=dbBack,colframe=dbFrame,colbacktitle=dbTitle,title={#1}}

\lstset{
  language=Python,basicstyle=\ttfamily\small,
  keywordstyle=\color{kbTitle}\bfseries,stringstyle=\color{wbTitle},
  commentstyle=\color{dbTitle},breaklines=true,frame=single,numbers=left,
  numberstyle=\tiny\color{gray},captionpos=b,columns=fullflexible,
  keepspaces=true,showstringspaces=false,extendedchars=true
}
\BeforeBeginEnvironment{lstlisting}{\Needspace{12\baselineskip}}

\hypersetup{colorlinks=true,linkcolor=kbTitle,urlcolor=kbTitle,citecolor=wbTitle}
\setlist{itemsep=0.25em,topsep=0.4em}
\setlength{\parindent}{2em}
\setlength{\parskip}{0.25em}
\newcolumntype{L}[1]{>{\raggedright\arraybackslash}p{#1}}

\providecommand{\notetitle}{Notas del curso: [tema del curso]}
\providecommand{\noteauthors}{Elaboradas a partir de los materiales públicos de Stanford CS25}
\providecommand{\notedate}{Versión reescrita en 2026}
\providecommand{\videochannel}{Stanford Online}
\providecommand{\videopublishdate}{2022}
\providecommand{\videoduration}{[duración del video]}
\providecommand{\videourl}{}
\providecommand{\videocoverpath}{cover.jpg}
\providecommand{\coursesite}{https://web.stanford.edu/class/cs25/}
\providecommand{\repourl}{https://github.com/hqhq1025/ai-course-notes}
\providecommand{\skillversion}{wdkns/wdkns-skills@39f1a04}

\newcommand{\figsource}[1]{\par\vspace{-0.3em}{\footnotesize\color{gray}\emph{Fuente: #1}}\par}
\newcommand{\teachervoice}[1]{\begin{knowledgebox}{Nota de clase}#1\end{knowledgebox}}
\newcommand{\term}[1]{\textbf{#1}}

\pagestyle{fancy}
\fancyhf{}
\fancyfoot[C]{\thepage}
\fancyfoot[R]{\footnotesize\color{gray}\href{\repourl}{AI Course Notes}}
\renewcommand{\headrulewidth}{0pt}
\renewcommand{\footrulewidth}{0pt}

\newcommand{\makecscover}{
\begin{titlepage}
\centering
\vspace*{0.7cm}
{\Large Stanford CS25 · Transformers United\par}
\vspace{0.8cm}
{\huge\bfseries \notetitle\par}
\vspace{0.6cm}
{\large \noteauthors\par}
\vspace{0.25cm}
{\large \notedate\par}
\vspace{0.9cm}
\ifdefempty{\videocoverpath}{}{
  \includegraphics[width=0.86\textwidth,height=0.43\textheight,keepaspectratio]{\videocoverpath}\par
}
\vfill
\begin{tcolorbox}[width=0.92\textwidth,colback=black!2!white,colframe=black!45,boxrule=0.7pt,arc=2pt]
\textbf{Autor o canal del video}: \videochannel\par
\textbf{Fecha de publicación}: \videopublishdate\par
\textbf{Duración del video}: \videoduration\par
\textbf{Sitio del curso}: \href{\coursesite}{\nolinkurl{\coursesite}}\par
\ifdefempty{\videourl}{}{\textbf{Enlace del video}: \href{\videourl}{\nolinkurl{\videourl}}\par}
\textbf{Normas de elaboración}: \skillversion\ + las normas source-first / teacher-voice / visual-QA de este repositorio
\end{tcolorbox}
\end{titlepage}
\tableofcontents
\newpage
}


\renewcommand{\notetitle}{Lecture 07: De Self-Attention a Non-Parametric Transformer---que el modelo razone entre datapoints}
\renewcommand{\noteauthors}{Elaborado a partir de la conferencia de Aidan Gomez, Jannik Kossen y Neil Band, los subtítulos manuales oficiales y el artículo original de NPT}
\renewcommand{\notedate}{CS25 V1 · Versión reescrita source-first de 2026}
\renewcommand{\videochannel}{Stanford Online}
\renewcommand{\videopublishdate}{2022-07-16}
\renewcommand{\videoduration}{1:05:43}
\renewcommand{\videourl}{https://www.youtube.com/watch?v=zejXBg-2Vpk}
\renewcommand{\videocoverpath}{cover.jpg}

\newcommand{\lecturefigure}[3]{%
\begin{figure}[H]
  \centering
  \includegraphics[width=0.96\textwidth]{slides-images/#1}
  \caption{#2}
  \figsource{#3}
\end{figure}
}

\begin{document}
\makecscover

\section{Auditoría de fuentes y mapa del curso: la misma attention, dos ejes de cómputo}

Esta clase se compone de dos partes que encajan entre sí. Primero, Aidan Gomez dedica unos quince minutos a repasar el origen del Transformer y la intuición detrás de su entrenamiento; después, Jannik Kossen y Neil Band presentan el Non-Parametric Transformer (NPT). Si solo se resumiera en el orden de la exposición, ambas partes parecerían «un repaso de fundamentos más un artículo»; sin embargo, el hilo didáctico real es más unitario: el Transformer coloca primero la attention \textbf{entre los tokens de una misma secuencia}, y NPT eleva después el eje sobre el que actúa la attention al plano \textbf{entre los datapoints de un mismo conjunto de datos}.

Esta nueva redacción toma como referencia temporal la publicación oficial de Stanford Online `zejXBg-2Vpk`. El video ofrece 1,476 subtítulos oficiales en inglés, hechos a mano y analizables, pero no hay un PDF de slides publicado por separado; por ello, se recuperaron de la grabación en 1080p y se seleccionaron manualmente 28 páginas didácticas. De las animaciones progresivas solo se conserva el estado final completo o los estados intermedios que tienen una función didáctica propia; se omiten las páginas anteriores a las que se regresa varias veces durante la sesión de Q\&A, las transiciones en blanco y el cierre del video, mientras que las explicaciones orales se incorporan al teacher-voice ledger y al cuerpo del texto.

\lecturefigure{slide-01-transformers-title.jpg}{Título de la primera parte: breve historia e intuición del Transformer.}{Video oficial de Stanford Online 00:00:46--00:01:07.}

Esta página de título recuerda que la tarea de Gomez no es volver a impartir un curso completo sobre el Transformer, sino extraer los tres mecanismos que determinan el resto de la historia: self-attention, multi-head attention y el autoregressive decoder que puede entrenarse en paralelo. Una vez comprendidos estos tres mecanismos, resulta más fácil ver que la innovación central de NPT no es «inventar otro tipo de attention», sino redefinir sobre qué objetos opera la attention.

\lecturefigure{slide-02-transformer-overview.jpg}{Panorama del Transformer: multi-head attention, self-attention y fast autoregressive decoding.}{Video oficial de Stanford Online 00:01:08--00:01:41.}

\begin{importantbox}{Una pregunta que condensa toda la clase}
Cuando el modelo predice un ejemplo, ¿solo puede apoyarse en «las características de ese ejemplo + la experiencia de entrenamiento ya comprimida en los parámetros», o puede leer directamente otros ejemplos de entrenamiento durante la inferencia? Una supervised neural network estándar suele adoptar la primera opción; NPT intenta aprender la segunda, y deja que la attention decida de extremo a extremo «qué ejemplos leer y cómo combinarlos».
\end{importantbox}

\teachervoice{Gomez resume la contribución del Transformer como tres componentes combinados, no como una fórmula aislada. Esta perspectiva es importante: LayerNorm, el learning-rate warmup, la inicialización y los detalles de las conexiones residual, que después se volvieron configuración estándar, eran en ese momento variables experimentales que podían determinar de forma notable el éxito o el fracaso.}

\subsection{Ruta de lectura}

A continuación se establece primero el marco matemático mínimo de la attention a nivel de token y del entrenamiento autoregresivo, y luego se pasa al input contract a nivel de conjunto de datos de NPT. La parte experimental no se limitará a repetir que «obtiene una muy buena posición en la clasificación», sino que responderá en orden a cuatro preguntas más exigentes: si NPT es suficientemente competitivo, si puede representar algoritmos que operan entre ejemplos, si de verdad utiliza otros ejemplos y qué estructura tiene esa dependencia.

\begin{knowledgebox}{No confundir los dos ejes}
\begin{tabularx}{\textwidth}{L{0.23\textwidth}>{\raggedright\arraybackslash}X >{\raggedright\arraybackslash}X}
\toprule
Mecanismo & Objetos que compara la attention & Pregunta principal \\
\midrule
Self-attention sobre secuencias & Los tokens / attributes dentro de un mismo ejemplo & ¿Qué elementos de la oración o del ejemplo están relacionados? \\
ABD de NPT & Los datapoints / rows de un mismo lote & ¿De qué otros ejemplos debe leer información el ejemplo actual? \\
ABA de NPT & Los attributes dentro de un mismo datapoint & ¿Cómo deben transformarse e interactuar los campos dentro de una fila? \\
\bottomrule
\end{tabularx}
\end{knowledgebox}

\subsection{Resumen de la sección}

Esta clase no es la unión de dos temas sin relación, sino la ampliación progresiva del ámbito sobre el que actúa la attention. La primera parte explica el entrenamiento en paralelo y las relaciones dentro de la secuencia; la segunda convierte el propio conjunto de datos en entrada del modelo, de modo que la predicción pueda depender explícitamente de otros datapoints.
\section{Self-Attention: de la alineación entre dos secuencias a las relaciones dentro de una secuencia}

La encoder--decoder attention tradicional suele hacer que la target-side query lea las source-side keys y values; la self-attention, en cambio, hace que query, key y value provengan del mismo conjunto de representaciones. Este cambio parece consistir solo en pasar el mismo tensor como los tres argumentos de la función, pero su significado estadístico cambia: el modelo ya no se limita a alinear secuencias distintas, sino que empieza a aprender un grafo de relaciones dentro de una misma muestra.

\lecturefigure{slide-03-self-attention.jpg}{La self-attention define source y target como la misma secuencia y aprende intra-sequence relations.}{Video oficial de Stanford Online, 00:01:42--00:02:39.}

Sea la representación de entrada $X\in\mathbb{R}^{m\times d_{\mathrm{model}}}$; la scaled dot-product attention se escribe como
\[
Q=XW_Q,\qquad K=XW_K,\qquad V=XW_V,
\]
\[
A=\operatorname{softmax}\!\left(\frac{QK^{\top}}{\sqrt{d_k}}\right),
\qquad Z=AV.
\]

\begin{itemize}
  \item $m$: longitud de la secuencia, es decir, el número de tokens que participan en el modelado de relaciones;
  \item $W_Q,W_K,W_V$: matrices aprendibles que proyectan la entrada en query, key y value;
  \item $d_k$: dimensión de cada key/query; $\sqrt{d_k}$ sirve para controlar la escala del dot product;
  \item $A_{ij}$: peso normalizado con el que el token $i$ lee información del token $j$;
  \item $Z_i$: nueva representación del token $i$, obtenida al agregar todos los values.
\end{itemize}

\teachervoice{El ponente usa «blue ball» para ilustrar la semántica de la self-attention: el modelo debería aprender a asociar el adjetivo blue con el sustantivo ball al que modifica. Este ejemplo no afirma que cierta head corresponda necesariamente a un árbol sintáctico; subraya que los attention weights pueden expresar hipótesis sobre las relaciones internas de la entrada.}

\begin{warningbox}{Los attention weights no son una explicación que se obtenga automáticamente}
$A_{ij}$ es el peso con el que se mezcla la información en una pasada hacia adelante. Puede indicar dónde está leyendo el modelo, pero por sí solo no demuestra una relación lingüística, una relación causal ni el motivo de una decisión. Para interpretarlo hay que considerar también el contenido de los values, las capas posteriores, la ablation o la intervention.
\end{warningbox}

\subsection{Multi-Head Attention: aprendizaje en paralelo de distintos subespacios de relaciones}

Un solo attention map solo puede comparar query y key en un único espacio de proyección. La multi-head attention proyecta la representación en varios subespacios de menor dimensión, lo que permite que distintas heads formen distintas medidas de similitud; después concatena los resultados y los vuelve a mezclar. Al leer la figura, conviene seguir primero el flujo de datos «división--attention independiente--concatenación», en lugar de entender el número de heads como una simple copia del modelo.

\lecturefigure{slide-04-multi-head-attention.jpg}{Multi-head attention: varios subespacios de relaciones se calculan de forma independiente y luego se concatenan.}{Video oficial de Stanford Online, 00:02:40--00:03:27.}

\[
\operatorname{head}_r
=\operatorname{Att}\!\left(QW_r^Q,KW_r^K,VW_r^V\right),
\]
\[
\operatorname{MHA}(Q,K,V)
=\operatorname{Concat}(\operatorname{head}_1,\ldots,\operatorname{head}_H)W^O.
\]

\begin{itemize}
  \item $H$: número de attention heads;
  \item $W_r^Q,W_r^K,W_r^V$: proyecciones independientes de la head $r$;
  \item $W^O$: vuelve a mezclar las salidas de todas las heads en la model dimension;
  \item cada head puede aprender relaciones distintas, pero ningún mecanismo garantiza una división del trabajo perfecta y nombrable por una persona.
\end{itemize}

\teachervoice{La intuición de Gomez es que una head puede inclinarse hacia la relación adjetivo--sustantivo y otra puede aprender otros patrones. Esto debe leerse como «la capacidad permite repartir las relaciones», no como «después del entrenamiento cada head se convierte automáticamente en un concepto humano claro».}

\begin{knowledgebox}{Por qué se divide entre $\sqrt{d_k}$}
Si las dimensiones de query y key son aproximadamente independientes y de varianza constante, la varianza del dot product crece con $d_k$. Los logits demasiado grandes hacen que la softmax se sature prematuramente y que los gradientes se reduzcan; el factor de escala devuelve los valores a un intervalo más estable.
\end{knowledgebox}

\subsection{El Transformer no es un único truco}

La self-attention y la multi-head attention son los títulos fáciles de recordar, pero Gomez enfatiza que el sistema original también dependía de un conjunto de decisiones de entrenamiento que en su momento no eran evidentes: LayerNorm, residual path, learning-rate warmup, la inicialización y los valores predeterminados de ingeniería de Tensor2Tensor. Tras la publicación del artículo, estos elementos se volvieron conocimiento común, lo que puede llevar a los lectores posteriores a creer que estos componentes, por naturaleza, debían combinarse así.

\teachervoice{El equipo realizó ablations repetidamente y encontró que retirar el warmup o la LayerNorm, o modificar configuraciones aparentemente menores, deterioraba de verdad la optimización. La conclusión de ingeniería es que el éxito de una arquitectura suele provenir de un conjunto de mecanismos que funcionan en conjunto, y que no se puede reducir toda la contribución a una sola fórmula de attention.}

\teachervoice{Gomez recuerda que el trabajo original tomó forma rápidamente en unos tres meses y se terminó a toda prisa antes de la fecha límite; en ese momento no previó su impacto. La lectura más valiosa de esta historia no consiste en fabricar un «momento de genialidad», sino en reconocer que las herramientas maduras, la iteración entre varias personas, la retroalimentación experimental y la oportunidad del momento contribuyeron juntas a un nuevo paradigma.}

\subsection{Resumen de la sección}

La self-attention traslada el modelado de relaciones de entre dos secuencias al interior de una misma entrada; la multi-head attention aporta varios subespacios de relaciones. La capacidad real de entrenar el Transformer depende además de un conjunto de técnicas de optimización y de conexiones residuales, y no puede explicarse solo con la expresión $QK^{\top}V$.
\section{Fast Autoregressive Training: por qué el decoder puede entrenarse en paralelo}

En la inferencia, un modelo generativo tiene que desplegar la salida token por token, pero el entrenamiento no necesita reproducir ese proceso en serie. Lo esencial del Transformer decoder es desplazar a la derecha toda la target sequence real antes de pasarla al modelo, y usar una causal mask para garantizar que cada posición solo vea el pasado. Así, la next-token loss de todas las posiciones se calcula en una sola pasada hacia adelante en paralelo.

\lecturefigure{slide-05-decoder-transition.jpg}{Transición entre el Transformer decoder y el fast autoregressive decoding.}{Video oficial de Stanford Online 00:03:28--00:04:07.}

Dada la entrada $x$ y la secuencia de salida $y_{1:T}$, la autoregressive factorization es
\[
p(y_{1:T}\mid x)=\prod_{t=1}^{T}p(y_t\mid y_{<t},x).
\]

\begin{itemize}
  \item $T$: longitud de la salida;
  \item $y_{<t}$: el prefix anterior a la posición $t$;
  \item en la inferencia hay que obtener primero $y_{t-1}$ para poder construir la condición del paso siguiente;
  \item en el entrenamiento ya se tiene la gold sequence completa, así que todos los prefixes correctos pueden construirse a la vez.
\end{itemize}

\lecturefigure{slide-06-sequence-generation.jpg}{Generación de secuencias: el objetivo del entrenamiento es aprender la correspondencia condicional entre la secuencia de entrada y la secuencia de salida.}{Video oficial de Stanford Online 00:04:08--00:05:19.}

\subsection{Teacher Forcing: entrenar cada posición con el gold prefix}

Teacher forcing (forzado por el maestro) consiste en darle al decoder, durante el entrenamiento, los tokens históricos reales y no las muestras que generó el propio modelo. Si la entrada desplazada a la derecha es $\tilde y=(\langle\mathrm{BOS}\rangle,y_1,\ldots,y_{T-1})$, una sola forward produce los logits de todas las posiciones y permite calcular
\[
\mathcal{L}_{\mathrm{TF}}
=-\sum_{t=1}^{T}\log p_{\theta}(y_t\mid y_{<t},x).
\]

\begin{itemize}
  \item $\theta$: parámetros del modelo;
  \item $p_{\theta}(y_t\mid y_{<t},x)$: probabilidad del $t$-ésimo target dado el prefix correcto;
  \item el hidden state de cada $t$ puede calcularse en paralelo, pero la dimensión de la secuencia sigue teniendo costo de attention y de memory;
  \item las condiciones del entrenamiento difieren de las de la generación libre, lo que produce exposure bias, pero eso no afecta la corrección del entrenamiento en paralelo.
\end{itemize}

\lecturefigure{slide-07-autoregressive-decoding.jpg}{La inferencia genera paso a paso; el entrenamiento, en cambio, puede recibir de una vez la target sequence correcta completa.}{Video oficial de Stanford Online 00:05:20--00:06:03.}

\teachervoice{Gomez dice que este detalle es «algo subtle, pero con un impacto enorme en la velocidad de entrenamiento». Aunque al principio del entrenamiento el modelo genere tokens sin sentido, sigue aprendiendo cada next-token conditional sobre el prefix correcto; precisamente esta supervisión en paralelo es lo que hace posible entrenar Transformers a gran escala.}

\lecturefigure{slide-08-teacher-forcing.jpg}{El problema central de teacher forcing: ¿cómo predecir todas las posiciones en una sola llamada al decoder sin filtrar la respuesta?}{Video oficial de Stanford Online 00:06:04--00:06:33.}

\begin{lstlisting}[caption={Pasos conceptuales de entrenamiento con teacher forcing y causal mask.}]
def autoregressive_training(source, target, model):
    decoder_input = shift_right(target, bos_token=True)
    causal_mask = upper_triangle_is_blocked(len(target))
    logits = model(source, decoder_input, causal_mask=causal_mask)
    return cross_entropy(logits, target)
\end{lstlisting}

\subsection{Causal Mask: paralelizar no significa permitir ver el futuro}

La subsección anterior explicó por qué la gold sequence permite obtener supervisión en paralelo; esta subsección añade la condición necesaria para que la técnica funcione: el grafo de cómputo de cada posición debe seguir siendo coherente con la autoregressive factorization. Si se entregara el target completo directamente a la self-attention, la posición $t$ podría leer $y_t$ o tokens futuros, y la tarea de entrenamiento se reduciría a copiar la respuesta; por eso, al leer la fórmula siguiente conviene revisar a la vez «qué posiciones son visibles» y «si se conserva la multiplicación de matrices en paralelo». La causal mask fija en menos infinito el triángulo superior estricto de los attention logits:
\[
C_{ij}=\begin{cases}
0, & j\le i,\\
-\infty, & j>i,
\end{cases}
\qquad
A=\operatorname{softmax}\!\left(\frac{QK^{\top}}{\sqrt{d_k}}+C\right).
\]

\begin{itemize}
  \item $i$: posición de la query actual;
  \item $j$: posición de key/value que se lee;
  \item las posiciones futuras con $j>i$ quedan con peso cero después de la softmax;
  \item la mask conserva el cálculo matricial en paralelo y, al mismo tiempo, mantiene la frontera de información de la autoregressive conditional.
\end{itemize}

\lecturefigure{slide-09-causal-mask.jpg}{Causal mask completa: cada posición solo puede leer su propio valor y el prefix a su izquierda.}{Video oficial de Stanford Online 00:12:36--00:13:04.}

\begin{knowledgebox}{Lectura de la figura: qué representa la matriz en blanco y negro}
El eje horizontal puede entenderse como las posiciones de key/value y el eje vertical como las posiciones de query. La región permitida forma un triángulo inferior; la región futura enmascarada no participa en la normalización. El valor de la mask en la figura no está en ahorrar cómputo, sino en impedir la target leakage, de modo que el entrenamiento en paralelo siga siendo equivalente a la next-token likelihood correcta.
\end{knowledgebox}

\begin{warningbox}{No confundir el paralelismo del entrenamiento con la decodificación en inferencia}
Teacher forcing hace que todas las posiciones se procesen en paralelo durante el entrenamiento, pero la autoregressive inference ordinaria sigue requiriendo un muestreo paso a paso, porque la entrada del paso siguiente depende del token que el modelo acaba de generar. Técnicas como KV cache y speculative decoding optimizan el costo de la inferencia, pero no forman parte de teacher forcing.
\end{warningbox}

\teachervoice{El ponente insiste en que «no se puede espiar hacia adelante»: ya pusimos la respuesta correcta en la decoder input; sin causal mask, un error de cero no significa que el modelo haya aprendido a generar, solo que descubrió una fuga de datos.}

\subsection{Resumen de la sección}

La autoregressive likelihood se sigue descomponiendo en el tiempo; teacher forcing solo aprovecha los targets conocidos para construir en paralelo todos los prefixes correctos. La causal mask es la condición necesaria para mantener la frontera de información: logra a la vez el «entrenamiento en paralelo» y que «el futuro no sea visible».
\section{El cambio de paradigma de NPT: leer explícitamente los datos de entrenamiento al predecir}

Ahora el objeto de la attention pasa de los tokens a los datapoints. Un parametric model estándar comprime la experiencia en los parámetros $\theta$ tras el entrenamiento y, al hacer inferencia sobre una muestra, por lo general solo lee esa muestra; NPT, en cambio, hace que los puntos de entrenamiento y los de prueba entren juntos al modelo, de modo que la predicción mantiene una dependencia explícita del conjunto de datos visible en ese momento.

\lecturefigure{slide-10-npt-overview.jpg}{Vista general de NPT en una diapositiva: dataset-as-input, datapoint self-attention y stochastic masking.}{Video oficial de Stanford Online, 00:17:00--00:17:21.}

Esta figura presenta tres componentes que no pueden suprimirse: primero, la entrada no es una sola fila, sino un conjunto de datos; segundo, el modelo aplica self-attention entre rows; tercero, el masking objective determina qué valores son visibles y cuáles deben reconstruirse. Si se conservan solo dos de ellos, cualesquiera que sean, no se obtiene lo que el artículo denomina el NPT predictive mechanism.

\lecturefigure{slide-11-npt-motivation.jpg}{Parametric prediction frente a NPT: en NPT, la fila objetivo lee explícitamente otros datapoints.}{Video oficial de Stanford Online, 00:17:22--00:18:17.}

\begin{importantbox}{Definición precisa de «non-parametric» en esta clase}
Non-parametric no equivale a «sin parámetros». NPT conserva una gran cantidad de parámetros aprendibles de attention, embedding y feed-forward. Lo que se subraya aquí es que la función de predicción depende explícitamente, durante la inference, de los supplied training data, lo que se escribe $p(y_*\mid x_*,\mathcal{D}_{\mathrm{train}})$, en lugar de depender solo de $x_*$ y de unos parámetros fijos $\theta$.
\end{importantbox}

\teachervoice{Alguien en redes sociales llamó a NPT «k-NN 2.0». La expresión sirve para formar una intuición: el modelo busca información en otras rows; pero NPT no se limita a buscar vecinos según una distancia fija, sino que puede aprender un proceso de match--transform--aggregate multicapa, dependiente de la tarea y no lineal.}

\subsection{Dónde quedan los non-parametric models tradicionales}

NPT no inventó la non-parametric prediction. Los Gaussian process, los kernel method y los k-nearest neighbors hacen que la predicción dependa explícitamente de los datos de entrenamiento; trabajos como deep kernel learning y neural process también intentan incorporar el representation learning. La tesis de NPT es más acotada: el Transformer puede funcionar como un aprendiz general de relaciones que entrena de extremo a extremo, de forma conjunta, la recuperación, la comparación y la agregación.

\lecturefigure{slide-12-traditional-nonparametric.jpg}{Non-parametric models tradicionales y sus extensiones profundas: el lugar histórico de NPT.}{Video oficial de Stanford Online, 00:22:00--00:23:39.}

\begin{knowledgebox}{Términos clave: tipos de explicit-data prediction}
\begin{tabularx}{\textwidth}{L{0.20\textwidth}L{0.27\textwidth}X}
\toprule
Método & Cómo usa los datos de entrenamiento & Relación con NPT \\
\midrule
k-NN & Recupera vecinos según una distancia predefinida y vota o promedia & NPT puede aprender la similitud, la agregación y las transformaciones posteriores; no se limita a una distancia fija. \\
Kernel method & Pondera los puntos de entrenamiento mediante kernel similarity & La attention también produce pesos que dependen de los datos, pero con representaciones profundas y composición en varias capas. \\
Gaussian process & El kernel define una distribución sobre funciones y proporciona la incertidumbre de la predicción & NPT no es un GP estándar ni posee automáticamente una Bayesian uncertainty coherente. \\
Deep kernel / neural process & Usan redes neuronales para aprender representaciones o una agregación condicional & Comparten con NPT la motivación de «aprendizaje parametrizado de relaciones + context explícito». \\
\bottomrule
\end{tabularx}
\end{knowledgebox}

\teachervoice{Kossen enfatiza que buscaban un non-parametric mechanism fácil de usar, plug-and-play y adaptable a múltiples escenarios, en lugar de depender de un inference scheme especializado y complejo. Este objetivo es una motivación de diseño; no debe interpretarse como prueba de que NPT ya supera a los métodos clásicos en todos los escenarios.}

\begin{warningbox}{NPT no es un sistema de meta-learning ya terminado en esta clase}
Los experimentos del artículo entrenan principalmente un modelo para un conjunto de datos fijo y suponen que train/test provienen de la misma distribución de tareas. El meta-learning suele exigir transferencia entre varias tareas o conjuntos de datos. NPT podría extenderse a un sistema que tome un nuevo dataset como context, pero esa es una línea futura que comentan los ponentes, no una capacidad que los experimentos de esta clase hayan demostrado.
\end{warningbox}

\teachervoice{En la sesión de Q\&A, los tres ponentes trazaron el límite con claridad: el trabajo actual sigue siendo supervised learning estándar; el meta-learning con múltiples conjuntos de datos, zero-shot o pocos gradient step es solo un framing que vale la pena explorar.}

\subsection{Resumen de la sección}

La novedad de NPT no es carecer de parámetros, sino conservar los datapoints de entrenamiento dentro del grafo de cómputo de la predicción y dejar que la attention aprenda a leerlos. Hereda la non-parametric intuition clásica y, al mismo tiempo, amplía el alcance del mecanismo aprendible mediante representaciones profundas y cómputo de relaciones en varias capas.
\section{Dataset-as-Input: matriz de datos, mask y contrato de predicción}

Para que todo el conjunto de datos entre al modelo, primero hay que reescribir la notación de aprendizaje supervisado (supervised-learning notation) como una matriz unificada. NPT trata las filas (rows) como datapoints y las columnas (columns) como atributos (attributes) con semántica compartida; la etiqueta (label) también ocupa una columna de atributo. La mask matrix determina qué entradas (entries) se observan y cuáles son incógnitas que el modelo debe predecir.

\lecturefigure{slide-13-entire-dataset-input.jpg}{Primer elemento: tomar el entire dataset, es decir, todos los datapoints, como entrada del modelo.}{Video oficial de Stanford Online 00:24:00--00:24:13.}

Sea
\[
X\in\mathbb{R}^{n\times d},\qquad M\in\{0,1\}^{n\times d}.
\]

\begin{itemize}
  \item $n$: número de datapoints en el forward actual; con datos pequeños puede ser el conjunto de datos completo, con datos grandes suele ser un minilote (mini-batch);
  \item $d$: número de atributos, que incluye las input features y puede incluir también las columnas objetivo (target columns);
  \item $X_{ij}$: el $j$-ésimo atributo del $i$-ésimo datapoint;
  \item $M_{ij}=1$: la entrada está enmascarada (mask) y debe predecirse; $M_{ij}=0$: la entrada es visible.
\end{itemize}

\lecturefigure{slide-14-datapoint-attention.jpg}{Segundo elemento: aplicar self-attention entre datapoints.}{Video oficial de Stanford Online 00:24:14--00:24:27.}

En un lote en tiempo de prueba (test-time batch), la attention puede conectar a la vez filas train--train, test--test y train--test. Aquí «train/test» es un papel determinado por la visibilidad del objetivo (target visibility): las filas de entrenamiento pueden exponer parte de las etiquetas, mientras que las etiquetas de las filas de prueba deben ocultarse por completo. El modelo ve la misma matriz, pero sabe por la mask qué entradas pueden servir como contexto.

\lecturefigure{slide-15-masking-objective-overview.jpg}{Tercer elemento: usar un masking-based objective para aprender de forma conjunta la reconstrucción del target y de las features.}{Video oficial de Stanford Online 00:24:28--00:24:41.}

Se definen las entradas masked y observed:
\[
X^{\mathcal{M}}=\{X_{ij}\mid M_{ij}=1\},\qquad
X^{\mathcal{O}}=\{X_{ij}\mid M_{ij}=0\}.
\]
NPT aprende
\[
p\!\left(X^{\mathcal{M}}\mid X^{\mathcal{O}},M\right),
\]
y produce $\widehat X\in\mathbb{R}^{n\times d}$. Esta representación convierte la clasificación, la regresión, la imputation, la self-supervision y la semi-supervision en «cambiar la posición de la mask».

\begin{importantbox}{La mask es la interfaz de la tarea}
Una misma architecture puede expresar problemas distintos mediante $M$: ocultar la columna de etiquetas de las filas de prueba produce supervised prediction; ocultar al azar entradas de features produce imputation/self-supervision; ocultar varias columnas objetivo produce multi-target prediction. Que la interfaz del modelo esté unificada no significa que la distribución de los datos y la loss de las distintas tareas queden unificadas automáticamente.
\end{importantbox}

\subsection{Notación completa: por qué también hay que dar la mask como entrada al modelo}

Las tres figuras anteriores presentaron el dataset, la datapoint attention y el masking objective; esta subsección los reúne en una interfaz de entrada sin ambigüedades. Al leer la figura de notación completa, conviene distinguir primero «cuál es el valor» de «si el valor es visible»: un mismo cero puede ser una observación válida o un marcador de posición para un dato faltante, y con solo mirar $X$ no es posible distinguirlos. Incrustar $M$ junto con $X$ equivale a indicarle explícitamente al modelo el epistemic status de cada valor: si es evidence conocida o una query por predecir; esto también determina en qué entradas debe aplicarse la loss.

\lecturefigure{slide-16-full-dataset-mask-notation.jpg}{Contrato de entrada completo: dataset $X$, mask $M$ y objetivo $p(X^{\mathcal{M}}\mid X^{\mathcal{O}})$.}{Video oficial de Stanford Online 00:25:06--00:27:15.}

\begin{knowledgebox}{Lectura de la figura: los tres colores y los signos de interrogación}
Cada fila es un datapoint y cada columna es un atributo de semántica fija. El signo de interrogación no significa «valor igual a cero», sino una entrada desconocida con $M_{ij}=1$. El modelo debe aprovechar a la vez las demás features de la misma fila, las features de otras filas y los targets de otras filas cuya exposición está permitida para reconstruir las posiciones marcadas con signo de interrogación.
\end{knowledgebox}

\begin{warningbox}{El límite infranqueable del test-label leakage}
El stochastic target masking puede conservar parte de las etiquetas en las filas de entrenamiento, lo que ayuda al modelo a aprender el lookup entre filas; las etiquetas reales de las filas de prueba nunca pueden usarse como entrada. Si al construir el lote o la mask se expone una test label, un resultado vistoso solo demuestra que el pipeline filtra la respuesta.
\end{warningbox}

\teachervoice{El ponente enuncia primero los tres elementos en conjunto y luego los desarrolla uno por uno: el conjunto de datos como entrada, la datapoint attention y el masking objective. Este orden nos recuerda que no conviene fijarse solo en el architecture block; lo que realmente define a NPT también incluye la input semantics y la training task.}

\subsection{Resumen de la sección}

NPT usa $(X,M)$ para representar de manera unificada la evidence y la query. Las filas son datapoints, las columnas son atributos de semántica compartida y la mask determina la tarea de predicción. Este contrato permite que los puntos train/test entren juntos a la attention y, al mismo tiempo, exige proteger estrictamente los test targets.
\section{NPT Architecture: attention alternada entre datapoints y attributes}

Una vez establecido el contrato de entrada, el siguiente paso es decidir a lo largo de qué dimensión opera la attention. NPT primero aplica a cada attribute un embedding sensible al tipo y obtiene un tensor tridimensional; después hace reshape repetidamente entre la «representación de la fila completa» y los «attributes dentro de la fila», y ejecuta de forma alternada Attention Between Datapoints (ABD) y Attention Between Attributes (ABA).

\subsection{Per-Attribute Embedding}

Las distintas columns pueden ser continuous, categorical, position/index o target. NPT no las une a la fuerza en un mismo tipo de token: para cada attribute usa el linear/typed embedding correspondiente y codifica junto con él la mask information.

\lecturefigure{slide-17-per-attribute-embedding.jpg}{Cada attribute se embebe de forma independiente y se forma un tensor de representación de $n\times d\times e$.}{Video oficial de Stanford Online 00:27:16--00:27:43.}

Se denota la representación inicial como
\[
H^{(0)}=\operatorname{Embed}(X,M)
\in\mathbb{R}^{n\times d\times e},
\]
donde $e$ es la dimensión del embedding de cada attribute. Para un categorical value puede usarse un learned lookup; para un continuous value, una proyección lineal con parámetros column-specific; el position/type embedding ayuda al modelo a conservar la column identity.

\begin{knowledgebox}{Por qué las rows son intercambiables y las columns no pueden intercambiarse libremente}
El orden de los datapoints normalmente no tiene significado, por lo que el modelo debe ser equivariant ante una row permutation. Las columns, en cambio, representan variables distintas, como edad, ingreso o etiqueta; intercambiar columns cambia la tarea, a menos que se transformen junto con ellas los column-specific embeddings, la mask y la semántica de la salida.
\end{knowledgebox}

\subsection{ABD: comprimir una fila completa en un attention token}

ABD necesita comparar datapoints, así que primero aplica flatten a los $d$ attribute embeddings de cada fila para obtener un vector de longitud $h=d\cdot e$. Después, la self-attention opera a lo largo de la dimensión $n$, y cada row puede agregar información de las demás rows.

\lecturefigure{slide-18-flatten-datapoint-attention.jpg}{ABD: tras aplicar flatten a los attributes, se hace self-attention entre los $n$ datapoints.}{Video oficial de Stanford Online 00:27:44--00:29:35.}

\[
R_{\mathrm{row}}:
\mathbb{R}^{n\times d\times e}\rightarrow
\mathbb{R}^{n\times (de)},
\]
\[
\operatorname{ABD}(H)
=R_{\mathrm{row}}^{-1}\!\left(
\operatorname{MHSA}(R_{\mathrm{row}}(H))
\right).
\]

\begin{itemize}
  \item $R_{\mathrm{row}}$: concatena los attributes de cada datapoint en una row representation;
  \item la sequence length de MHSA es aquí $n$, no $d$;
  \item una ABD layer aprende pairwise row interactions, y varias capas pueden combinarse en higher-order interactions;
  \item el término dominante de cómputo y memory crece con $n^2$.
\end{itemize}

\subsection{ABA: reorganizar los attributes dentro de cada fila}

Si solo se usa ABD, la fila completa se trata como un token indivisible y resulta difícil recodificar columns concretas. ABA devuelve el tensor a la forma $n\times d\times e$ y, para cada fila de manera independiente, hace self-attention a lo largo de los $d$ attributes, de modo que el modelo aprende una per-datapoint transformation.

\[
\operatorname{ABA}(H)
=\operatorname{stack}_{i=1}^{n}
\left[\operatorname{MHSA}(H_{i,:,:})\right]
\in\mathbb{R}^{n\times d\times e}.
\]

\begin{itemize}
  \item $H_{i,:,:}$: la attribute sequence de la fila $i$;
  \item ABA comparte parámetros entre las rows, pero cada fila se calcula de forma independiente;
  \item aprende attribute interactions y no lee directamente otros datapoints;
  \item la alternancia entre ABD y ABA hace que «encontrar filas relacionadas» y «transformar la fila actual» se refuercen mutuamente.
\end{itemize}

\lecturefigure{slide-19-abd-aba-architecture.jpg}{Arquitectura alternada completa: ABD aprende relaciones entre datapoints, ABA aprende transformaciones de una sola fila, y se conserva la row permutation equivariance.}{Video oficial de Stanford Online 00:29:36--00:31:03.}

\begin{importantbox}{Permutation equivariance}
Sea $P\in\{0,1\}^{n\times n}$ una row permutation matrix; NPT satisface
\[
f(PX,PM)=P\,f(X,M).
\]
Es decir, permutar las rows de entrada solo permuta de la misma manera las rows de salida, y no debe cambiar el contenido de la predicción de cada datapoint. Equivariance no es lo mismo que invariance: la salida no permanece idéntica, sino que sigue la permutación de la entrada.
\end{importantbox}

\teachervoice{El ponente enfatiza que ABD y ABA tienen funciones distintas: el primero captura higher-order relationships between datapoints y el segundo se ocupa de las individual datapoint transformations. Llamar a ambos «attention block» hace perder el tensor axis más importante.}

\teachervoice{En la sesión de Q\&A, los autores reconocen que ABA no es un dogma insustituible: la ablation muestra que es útil, pero si el número de attributes es muy grande, puede considerarse reemplazarlo por un compact MLP para reducir el costo de escalar en la attribute dimension.}

\subsection{Sistema de tres etapas y masking objective}

Las dos primeras etapas resuelven «cómo se representa la entrada y cómo se propagan las relaciones»; la tercera determina por qué el modelo aprende a hacer lookup. NPT no entrena primero un classifier ordinario para luego agregar neighbors de forma improvisada durante la inferencia: desde el inicio del entrenamiento se enfrenta una y otra vez a features y targets ocultos, y debe reconstruirlos a partir de las entries visibles.

\lecturefigure{slide-20-three-stage-overview.jpg}{Las tres etapas de NPT: dataset input, datapoint attention y masking-based objective.}{Video oficial de Stanford Online 00:31:04--00:33:33.}

\lecturefigure{slide-21-stochastic-masking-objective.jpg}{Objetivo conjunto de feature masking y target masking.}{Video oficial de Stanford Online 00:33:34--00:33:52.}

Sea $\mathcal{L}_{\mathrm{Targets}}$ la negative log-likelihood o la loss de regresión de los targets con mask, y $\mathcal{L}_{\mathrm{Features}}$ la loss auxiliar de reconstrucción de las feature entries con mask aleatoria:
\[
\mathcal{L}_{\mathrm{NPT}}
=(1-\lambda)\mathcal{L}_{\mathrm{Targets}}
+\lambda\mathcal{L}_{\mathrm{Features}}.
\]

\begin{itemize}
  \item $p_{\mathrm{feature}}$: probabilidad de ocultar aleatoriamente feature values durante el entrenamiento;
  \item $p_{\mathrm{target}}$: probabilidad de ocultar aleatoriamente training targets durante el entrenamiento;
  \item $\lambda$: equilibrio entre la target loss y la auxiliary feature loss;
  \item en test time solo se aplican mask y evaluación a los test targets, cuyo valor verdadero no puede exponerse.
\end{itemize}

\begin{knowledgebox}{Qué hace cada masking term}
El feature masking hace que el modelo aprenda la estructura de toda la matriz de datos, aumenta la supervision y aporta regularization; en la ablation del artículo, ayudó en ocho de diez tabular datasets. El target masking, por su parte, mantiene visibles parte de las training labels, de modo que otras rows ocultas puedan aprender a leer, emparejar y transformar esas labels.
\end{knowledgebox}

\teachervoice{La explicación central de Neil Band es que el modelo no necesita memorizar en sus parameters cada training input--output mapping; puede dedicar su capacidad a aprender «cómo aprovechar las demás training features y targets». Es justo aquí donde la learned k-NN intuition deja de ser una metáfora y se convierte en un mecanismo a nivel de objective.}

\begin{lstlisting}[caption={Pseudocódigo conceptual del forward de NPT y de la mask aleatoria.}]
def npt_step(dataset, train_rows, test_rows, model):
    mask = sample_feature_masks(dataset)
    mask |= sample_training_target_masks(train_rows)
    mask |= all_test_targets(test_rows)

    hidden = per_attribute_embed(dataset, mask)
    for _ in range(num_blocks):
        hidden = attention_between_datapoints(hidden)  # ABD
        hidden = attention_between_attributes(hidden)  # ABA
    prediction = project_all_attributes(hidden)
    return masked_entry_loss(prediction, dataset, mask)
\end{lstlisting}

\begin{warningbox}{La aproximación por minilotes cambia el retrieval set}
Al usar minilotes aleatorios con datos grandes, una query solo puede leer los datapoints del mismo lote, no todo el conjunto de entrenamiento. Aumentar el tamaño del lote mejora la cobertura de candidatos, pero implica un costo de attention de $O(b^2)$; por ello, el uso de minilotes es un compromiso estadístico y de sistemas, no un cálculo estrictamente equivalente sobre todos los datos.
\end{warningbox}

\teachervoice{Los autores dan una escala concreta: sin minilotes se pueden manejar alrededor de 8,000 points; un conjunto de datos de 11 million points necesariamente depende de minilotes. También observaron que no hace falta un lote «absurdamente grande» para aprender relaciones entre datapoints, aunque esto sigue siendo una conclusión empírica.}

\subsection{Resumen de la sección}

NPT embebe $(X,M)$ en un tensor de $n\times d\times e$, alterna ABD para procesar las rows y ABA para procesar las columns, y entrena un lookup relacional con un masking objective. La row equivariance, la test-label protection y la quadratic datapoint attention son las tres restricciones firmes para entender la implementación.
\section{Benchmark Evidence: por qué elegir tabular data y cómo leer las tablas}

Si solo se prueba en benchmarks que las redes profundas dominan con facilidad, es difícil saber si la dataset-level attention aporta valor real. Los autores eligen tabular data porque abarca clasificación y regresión, campos continuos y discretos, y desde cientos hasta decenas de millones de muestras, además de que el tree-based boosting ha sido fuerte durante mucho tiempo; esto obliga a comparar NPT con baselines no neuronales que de verdad compiten.

\lecturefigure{slide-22-tabular-domain.jpg}{Tabular benchmark: la escala, los tipos de campo, las tareas y los baselines fuertes son muy diversos.}{Video oficial de Stanford Online 00:37:16--00:38:07.}

\begin{knowledgebox}{Referencia rápida de baselines y métricas}
\begin{tabularx}{\textwidth}{L{0.19\textwidth}L{0.25\textwidth}>{\raggedright\arraybackslash}X}
\toprule
Nombre & Mecanismo & Por qué es necesario compararlo \\
\midrule
XGBoost / LightGBM & gradient-boosted decision trees eficientes & Baseline fuerte de la industria para tabular data; maneja bien las divisiones no lineales y las características faltantes o dispersas. \\
CatBoost & ordered boosting orientado a categorical features & Evita que NPT obtenga una ventaja falsa solo por tratar mejor los campos categóricos. \\
Random Forest / GBDT & Modelos de árboles con bagging o boosting & Cubre la ensemble family clásica. \\
MLP & parametric neural baseline que procesa cada fila de forma independiente & Verifica si la mejora proviene del mecanismo entre datapoints. \\
TabNet & attentive neural model orientado a tabular data & Compara con una domain-specific neural architecture. \\
k-NN & local lookup con distancia fija & Contrasta si el «learned non-parametric lookup» supera a los vecinos simples. \\
AUROC / Accuracy / RMSE & Ordenamiento en clasificación binaria, tasa de aciertos en clasificación multiclase, error de regresión & Las métricas de los tres tipos de tarea tienen direcciones distintas; el artículo las unifica convirtiéndolas en method rank. \\
\bottomrule
\end{tabularx}
\end{knowledgebox}

\teachervoice{Los autores no eligieron el tabular domain solo porque los datos fueran convenientes, sino porque las general-purpose deep nets suelen perder aquí frente al boosting; si NPT solo superara a un MLP débil, eso no bastaría para sostener su afirmación de generalidad.}

\teachervoice{Al equipo le sorprendió la robustness de NPT en small datasets, y eso sin realizar una cantidad extrema de ajustes de hyperparameters. Vale la pena registrar esta observación verbal, pero no sustituye una comparación rigurosa del tuning-budget.}

\subsection{Average rank no es pooled accuracy}

La página de resultados agrupa los diez UCI datasets por tarea: cuatro de binary classification usan AUROC, dos de multi-class classification usan accuracy y cuatro de regression usan RMSE. Para cada dataset primero se ordenan los methods y luego se promedia dentro del mismo grupo de tareas; por ello, un valor menor es mejor, y no se pueden comparar directamente columnas distintas como si fueran la misma magnitud física.

\lecturefigure{slide-23-tabular-benchmarks.jpg}{Rango promedio en el UCI benchmark y primeros resultados en CIFAR-10.}{Video oficial de Stanford Online 00:38:08--00:40:44.}

\begin{table}[H]
\centering
\small
\caption{Rangos promedio clave de la Table 1 del artículo; un valor menor es mejor.}
\begin{tabularx}{\textwidth}{L{0.22\textwidth}L{0.23\textwidth}L{0.23\textwidth}X}
\toprule
Método & Binary AUROC rank & Multi-class accuracy rank & Regression RMSE rank \\
\midrule
NPT & $2.50\pm0.87$ & $2.50\pm0.50$ & $3.25\pm1.31$ \\
CatBoost & $2.75\pm0.85$ & $3.50\pm0.50$ & $3.00\pm0.91$ \\
XGBoost & $4.75\pm1.25$ & $2.50\pm1.50$ & $3.25\pm0.63$ \\
MLP & $5.75\pm1.49$ & $3.00\pm2.00$ & $5.00\pm1.22$ \\
k-NN & $8.25\pm0.48$ & $8.50\pm0.50$ & $8.75\pm0.25$ \\
\bottomrule
\end{tabularx}
\end{table}

\begin{knowledgebox}{Lectura de la figura: primero el intervalo competitivo, luego el número de primeros lugares}
En binary y multi-class, el rank promedio de NPT está en el grupo de cabeza, y en regression queda cerca de CatBoost/XGBoost; el artículo también informa que NPT es top performer en cuatro de los diez tabular datasets. La conclusión más prudente es «competitivo en varios tipos de tarea», no «supera de manera uniforme a todos los tree methods».
\end{knowledgebox}

\begin{warningbox}{Qué oculta el rango promedio}
Diez datasets son pocos, y cada grupo de tareas tiene solo de dos a cuatro conjuntos de datos; el rank pierde el performance gap absoluto y además depende del tuning protocol y de la dataset selection. El error estándar proviene de la rank variation entre datasets; no es el intervalo de confianza de una sola accuracy medida sobre una muestra grande.
\end{warningbox}

\teachervoice{En la sesión de Q\&A se dedicó bastante tiempo a aclarar «4 of 10», «2 of 10» y el standard error. Por eso las notas deben dejar claro que estas cifras describen primero la agrupación por dataset/task y solo después el rango promedio de los methods; a partir de los decimales en apariencia pequeños no puede inferirse una ventaja significativa general.}

\subsection{Los límites correctos del image result}

El artículo también prueba NPT en CIFAR-10 y MNIST. El resultado en CIFAR-10 que se muestra en la clase usa un CNN encoder más NPT y reporta una accuracy de 93.7\%; esto indica que la datapoint interaction entre distintas images puede integrarse en un pipeline de visión, pero no demuestra que NPT sustituya al backbone visual. Los resultados con linear patching del apéndice del artículo son más débiles, lo que también refleja que el input embedding sigue influyendo en el desempeño.

\begin{warningbox}{«Formato de entrada general» no equivale a «el front-end es irrelevante»}
Que cualquier modality pueda hacer reshape a una matrix solo indica que la interfaz puede recibirla; si se conserva la estructura local, si existe un encoder adecuado y cómo se configuran el volumen de datos y la augmentation siguen determinando la optimization y la sample efficiency. La contribución de NPT es agregar un mecanismo entre datapoints, no eliminar todo el conocimiento propio de cada modality.
\end{warningbox}

\subsection{Resumen de la sección}

NPT entra en la zona competitiva de los baselines fuertes en diversas tabular tasks y muestra la viabilidad de la image extension. La evidencia respalda que es «competitivo», pero el rank promedio, el número limitado de datasets y los distintos front-ends no nos permiten afirmar que sea el mejor en todos los escenarios.
\section{Corruption Experiment: ¿el modelo realmente usa otros datapoints?}

El benchmark performance no demuestra directamente que la attention entre datapoints cumpla una función: un NPT de alta capacidad podría aprender a ignorar las demás rows y degenerar en un per-row parametric model común. Por ello, los autores diseñaron un intervention-like test que, mientras mantiene intacta la target row actual, destruye la información relacional de todas las demás rows y luego observa si la predicción empeora.

\lecturefigure{slide-24-corruption-goal.jpg}{Objetivo del corruption experiment: impedir que el modelo obtenga relaciones útiles de otros datapoints.}{Video oficial de Stanford Online 00:51:16--00:51:31.}

Para cada target entry evaluada $(i,d)$, el método conserva la fila $i$ y aplica a cada attribute column de las demás rows una permutation aleatoria independiente. Puede escribirse como
\[
\widetilde X_{k,j}=X_{\pi_j(k),j},
\qquad k\ne i,
\]
donde cada column $j$ usa una permutation independiente $\pi_j$.

\begin{itemize}
  \item La query row actual $i$ no cambia;
  \item los marginal values de cada column se conservan aproximadamente;
  \item las correspondencias feature--feature y feature--target que existían dentro de una misma row se rompen;
  \item si la predicción depende de la estructura de las demás rows, la performance debería bajar tras la corruption.
\end{itemize}

\lecturefigure{slide-25-corruption-method.jpg}{Corruption method: se permutan las demás rows attribute por attribute y se conserva la query row.}{Video oficial de Stanford Online 00:51:32--00:53:01.}

\teachervoice{El ponente aclaró expresamente que no quería «estropear» la entrada con ruido simple, porque el modelo podría limitarse a detectar un distribution shift. La permutation column por column conserva en lo posible las marginal statistics, pero destruye las correspondencias aprovechables entre campos y entre muestras.}

\subsection{Resultados: algunas tareas dependen fuertemente, otras ignoran activamente}

Si el NPT dependiera por igual de las demás rows en todos los conjuntos de datos, la corruption debería causar pérdidas grandes en general; los resultados reales son más matizados: tareas como Protein empeoran notablemente, mientras que Forest, Kick, Breast Cancer y otras casi no cambian. Esto indica que la architecture ofrece un relation channel, pero la optimization puede decidir, según los datos, si lo usa o no.

\lecturefigure{slide-26-corruption-results.jpg}{Performance change después de destruir la información de los demás datapoints.}{Video oficial de Stanford Online 00:53:02--00:54:43.}

\begin{table}[H]
\centering
\small
\caption{Table 2 del artículo: cambio de performance después de la corruption; los valores negativos indican empeoramiento.}
\begin{tabularx}{\textwidth}{L{0.21\textwidth}*{4}{>{\centering\arraybackslash}X}}
\toprule
Conjunto de datos de clasificación & CIFAR-10 & Poker & Income & Higgs \\
\midrule
$\Delta$ Accuracy & $-1.2$ & $-1.1$ & $-1.1$ & $-0.5$ \\
\bottomrule
\end{tabularx}
\vspace{0.6em}
\begin{tabularx}{\textwidth}{L{0.21\textwidth}*{4}{>{\centering\arraybackslash}X}}
\toprule
Conjunto de datos de clasificación & MNIST & Forest & Kick & Breast Cancer \\
\midrule
$\Delta$ Accuracy & $-0.4$ & $-0.1$ & $-0.1$ & $0.0$ \\
\bottomrule
\end{tabularx}
\vspace{0.6em}
\begin{tabularx}{\textwidth}{L{0.21\textwidth}*{4}{>{\centering\arraybackslash}X}}
\toprule
Conjunto de datos de regresión & Yacht & Protein & Boston & Concrete \\
\midrule
$\Delta\mathrm{RMSE}/\mathrm{RMSE}$ & $-52\%$ & $-21\%$ & $-20\%$ & $-7\%$ \\
\bottomrule
\end{tabularx}
\end{table}

\begin{importantbox}{Lo más interesante no es que «todo baje»}
En algunos datasets, el NPT descubre que las demás rows no aportan suficiente valor, de modo que puede aprender a ignorarlas y degenerar en un predictor aproximadamente parametric; en otros datasets, la corruption provoca un colapso grave de la performance. Lo que se aprende no es la adopción fija de una non-parametric rule, sino el grado de dependencia.
\end{importantbox}

\begin{warningbox}{La corruption respalda la dependencia del mecanismo, pero no es una prueba causal completa}
Un performance drop indica que la predicción aprovechó la structure entre rows que fue destruida; por sí solo, no nos dice cuál es la única ruta causal de cierta attention head, ni permite concluir que «lo non-parametric es en general superior a lo parametric». La corruption también altera la distribución conjunta de orden superior, por lo que todavía podría introducir algún out-of-distribution pattern al que el modelo sea sensible.
\end{warningbox}

\teachervoice{Tanto el artículo como la clase señalan con cautela que el hecho de que Forest, Kick y Breast Cancer casi no se vean afectados no es un fallo del experimento, sino que el modelo posiblemente juzgó que el lookup entre datapoints no ofrecía ventaja. Esta posibilidad de degenerar respalda justamente la afirmación de que «el grado de dependencia se aprende de los datos».}

\subsection{Resumen de la sección}

El corruption test avanza de «el modelo tiene attention entre filas» a «el modelo realmente depende de la información entre filas en varias tareas reales». La evidencia proviene de un performance change de tipo contrafactual, pero el experimento no puede identificar un mecanismo único ni generalizarse a una conclusión universal sobre la superioridad de lo parametric o lo non-parametric.
\section{Duplicate e Intervention: descomponer la learned lookup en pasos observables}

La corruption sobre datos reales muestra que el canal relacional se utiliza, pero aún no deja claro qué algoritmo aprendió el modelo. La semi-synthetic duplicate task construye un mecanismo más transparente: se copian varias rows, se ocultan los target de las originals y se exponen los target de las duplicates; una predicción correcta debe emparejar la row correspondiente, leer el target y luego copiarlo de vuelta a la query.

\lecturefigure{slide-27-duplicate-intervention.jpg}{Cadena completa de evidencia de la semi-synthetic duplicate task, la attention, las predicciones y la target intervention.}{Video oficial de Stanford Online, 00:54:44--01:00:41.}

\begin{knowledgebox}{Lectura de la figura: orden de razonamiento de los cinco panel}
\begin{enumerate}
  \item (a) Construir las original/duplicate rows y ocultar los targets de las originals;
  \item (b) verificar si la ABD attention se concentra en las duplicates correspondientes;
  \item (c) comprobar que las predictions tienen una correlación alta con los duplicate targets;
  \item (d) modificar deliberadamente un duplicate target en test time, sin volver a entrenar;
  \item (e) observar si la original prediction cambia de acuerdo con la intervention.
\end{enumerate}
\end{knowledgebox}

Esta tarea exige que el modelo aprenda un programa compuesto:
\[
x_i\xrightarrow{\text{match features}}
x_{j(i)}^{\mathrm{dup}}
\xrightarrow{\text{read target}}
y_{j(i)}^{\mathrm{dup}}
\xrightarrow{\text{copy/transform}}
\widehat y_i.
\]

\begin{itemize}
  \item $j(i)$: el duplicate index que corresponde a la original row $i$;
  \item el match lo realiza la learned attention representation, no un row ID escrito de antemano;
  \item la intervention modifica $y_{j(i)}^{\mathrm{dup}}$ y deja sin cambios el resto del mecanismo;
  \item si $\widehat y_i$ se desplaza de manera estable con la intervention, el modelo usa una lookup path generalizable.
\end{itemize}

\teachervoice{El ponente reporta que, después de la intervention, la correlación entre la prediction y el duplicate value se mantiene en aproximadamente 99.6\%. Lo decisivo es que el modelo, sin volver a entrenarse, sigue un target value que está fuera de la distribución de entrenamiento; esto se acerca más a una prueba del mecanismo que observar solo el test error habitual.}

\subsection{Por qué no es solo nearest neighbor}

Un plain k-NN también puede encontrar la duplicate y copiar el target, por lo que la tarea básica no basta para mostrar la transformación relacional aprendible de NPT. Los autores van más allá y le suman uno al duplicate target; una copy rule fija se equivoca de manera sistemática, mientras que NPT puede aprender a hacer primero la lookup y después restar uno. De forma más general, NPT puede aprender una joint feature--target relation, en lugar de solo promediar localmente bajo una metric fija.

\begin{importantbox}{Lo que aprende NPT no es una tabla estática de vecinos}
La evidencia más fuerte no es que «la attention se dirige a rows similares», sino que el modelo, después del match, todavía puede ejecutar una task-dependent transformation. Las múltiples capas ABD/ABA permiten que la lookup, el attribute reasoning y la value transformation se combinen en un learned algorithm.
\end{importantbox}

\teachervoice{Kossen reconoce por iniciativa propia que la duplicate task básica puede resolverse con nearest neighbor, y luego propone la transformación de sumar uno para aumentar la dificultad. Vale la pena aprender esta forma de argumentar: no evitar los baselines sencillos, sino construir el experimento mínimo que distinga la capacidad del mecanismo.}

\begin{warningbox}{Límites para extrapolar el semi-synthetic success}
La tarea se diseñó deliberadamente para que existan duplicates claras e interventions verificables, por lo que sirve para demostrar la representational capability. Los datos reales normalmente no tienen coincidencias perfectas, un mecanismo único y limpio ni labels manipulables; el éxito no equivale a que NPT ya domine el causal reasoning en general.
\end{warningbox}

\subsection{Resumen de la sección}

El duplicate/intervention experiment descompone el razonamiento entre datapoints en cuatro pasos: match, read, transform y write, y usa la test-time intervention para comprobar que la prediction cambia con el context. Es una evidencia fuerte del mecanismo, pero sigue limitada a una semi-synthetic construction transparente.
\section{Limitaciones, conexiones y futuro: cuál es el costo de la dataset attention}

La capacidad expresiva de NPT proviene de hacer que los datapoints sean visibles entre sí, y ese mismo diseño crea también el mayor cuello de botella del sistema. Si el lote contiene $b$ datapoints y la hidden dimension de cada fila, después del flatten, es $h=de$, el ABD attention map tiene tamaño $b\times b$, y el cómputo principal puede escribirse de forma aproximada como
\[
\operatorname{Cost}_{\mathrm{ABD}}=O(b^2h),
\qquad
\operatorname{Memory}_{\mathrm{attn}}=O(b^2).
\]

\begin{itemize}
  \item $b$: número de datapoints que participan efectivamente en una relational lookup;
  \item $h$: hidden width de cada row;
  \item aumentar $b$ amplía la cobertura de context candidato, pero con un quadratic cost;
  \item el mini-batch aleatorio reduce el costo, pero puede omitir neighbors potencialmente clave que quedan fuera del lote.
\end{itemize}

\lecturefigure{slide-28-summary-future-work.jpg}{Síntesis de NPT: experimentos, limitación de quadratic scaling y future work.}{Video oficial de Stanford Online, 01:00:42--01:05:36.}

\begin{knowledgebox}{Lectura de la figura: conclusiones y preguntas abiertas por separado}
A la izquierda están los resultados que esta clase ya mostró: dataset-as-input, datapoint self-attention, tabular/image benchmarks e intervention evidence. A la derecha están las limitaciones aún sin resolver: el scaling y tareas nuevas como multi-task, few-shot, semi-supervised y domain adaptation. El future work no es una lista de capacidades actuales.
\end{knowledgebox}

\subsection{Conexión con GNN}

Si cada datapoint se considera un node, ABD equivale a realizar learned message passing sobre un fully connected graph. La diferencia es que muchas GNN reciben sparse edges definidos de antemano, mientras que NPT aprende las relaciones de forma suave a partir del attention score; sin embargo, los graph models modernos también pueden aprender edges, de modo que la frontera entre ambos no es absoluta.

\teachervoice{En la sesión de Q\&A, el expositor acepta la analogía con el fully connected graph y señala que NPT puede entenderse como discovering relational structure. Neil Band menciona Neural Relational Inference: cuando los edges son desconocidos, tanto attention como message passing aprenden interacciones latentes.}

\begin{warningbox}{Un fully connected graph no implica que las relaciones se hayan descubierto correctamente}
Permitir la comunicación entre cualquier par de puntos solo define un hypothesis space. El modelo todavía puede aprender un shortcut, una spurious similarity o una attention casi uniform; para determinar si las relaciones son útiles se necesitan corruption, intervention, OOD evaluation y sparse retrieval escalable.
\end{warningbox}

\subsection{El valor distinto de small data y large data}

La ventaja de small data es que el conjunto de datos completo cabe en un solo forward, de modo que el modelo obtiene un context casi completo; large data, en cambio, nos obliga a introducir retrieval, sparsity, kernel approximation o learned representatives. Aunque estas opciones sacrifican la attention exacta con conexión total, pueden vincular NPT con el retrieval-augmented learning moderno: los parámetros se encargan de aprender las reglas de recuperación y transformación, y los datapoints externos se encargan de almacenar hechos actualizables.

\teachervoice{Al final, el expositor plantea una pregunta incisiva: ¿cuántos parámetros de los grandes parametric models se dedican en realidad a «almacenar datos»? Si una lookup explícita puede asumir parte de la memoria, los parámetros quizá podrían concentrarse más en aprender cómo recuperar, combinar y razonar. Esta pregunta cobró después todavía más importancia en los retrieval-augmented systems.}

\begin{knowledgebox}{Hoja de ruta de scaling}
\begin{tabularx}{\textwidth}{L{0.24\textwidth}>{\raggedright\arraybackslash}X >{\raggedright\arraybackslash}X}
\toprule
Línea & Beneficio & Riesgo nuevo \\
\midrule
Random mini-batch & Implementación sencilla y entrenamiento estable & Un datapoint clave puede no estar en el lote; el context es aleatorio. \\
Sparse / local attention & Reduce $b^2$ a una cantidad pequeña de edges & Requiere definir o aprender un candidate graph. \\
ANN / learned retrieval & Primero recupera las top-$k$ rows y luego aplica NPT reasoning & El recall del retriever se vuelve el cuello de botella previo. \\
Representative points & Comprime el conjunto de datos con prototypes / inducing points & La compresión puede perder muestras raras y detalles locales. \\
Kernel / low-rank approximation & Aproxima la attention global & El error de aproximación y la eficiencia en hardware deben evaluarse en conjunto. \\
\bottomrule
\end{tabularx}
\end{knowledgebox}

\subsection{Del conjunto de datos fijo hacia el meta-learning}

Los experimentos de esta clase corresponden a fixed-dataset supervised learning, pero el input contract se presta de forma natural a una context-based extension: se colocan el support set, el unlabeled set y el query set en una matriz unificada, y la tarea se especifica mediante masks. Si se entrena a través de varios datasets, NPT podría aprender un within-dataset algorithm general; esto lo conecta directamente con in-context learning, few-shot learning, semi-supervised learning y domain adaptation.

\teachervoice{El expositor clasifica explícitamente estas líneas como future work, entre ellas eliminar la mini-batch approximation, la few-shot generalization, domain adaptation, semi-supervised y multi-task learning. Estas notas las presentan solo como líneas de investigación y no describen experimentos aún no realizados como capacidades ya establecidas.}

\subsection{Resumen de la sección}

La capacidad expresiva de la dataset attention y su costo $O(b^2)$ provienen del mismo diseño. GNN, retrieval, sparse attention y meta-learning ofrecen extensiones naturales, pero cada aproximación cambia el context visible y el statistical contract, por lo que requiere validar de nuevo el mecanismo.
\section{Síntesis y ampliación}

\subsection{De token relation a datapoint relation}

La visión unificada que más vale conservar de esta clase es la siguiente: la attention no pertenece por naturaleza a los token del lenguaje; es una plantilla de cómputo que «permite a un objeto actualizar su representación a partir de un conjunto de otros objetos». El Transformer aprende relaciones entre tokens; NPT convierte cada datapoint en un objeto y, mediante reshape, razona alternando entre dos ejes: datapoints y attributes.

\begin{importantbox}{Doce conclusiones centrales}
\begin{enumerate}
  \item La self-attention hace que query, key y value provengan del mismo conjunto, y aprende las relaciones internas de la entrada.
  \item La multi-head attention ofrece varios subespacios de proyección, pero no garantiza que cada head pueda recibir un nombre interpretable por una persona.
  \item Teacher forcing usa gold prefixes para entrenar en paralelo todas las next-token positions.
  \item La causal mask impide que el modelo lea los targets futuros; es el límite que asegura la corrección de la likelihood.
  \item NPT sigue siendo una parametric neural network; «non-parametric» significa que la inference depende explícitamente de los datapoints de entrenamiento.
  \item $(X,M)$ coloca la observed evidence y las masked queries en una misma dataset representation.
  \item ABD aplica attention entre rows; ABA aplica attention entre los attributes de cada fila.
  \item La row permutation equivariance garantiza que el orden de los datapoints no aporte semántica espuria.
  \item El feature masking proporciona regularization; el target masking entrena el relational lookup.
  \item La evidencia de UCI average-rank respalda que el modelo es competitivo, pero no que supere a los demás en todas las tareas.
  \item Las intervenciones de corruption y duplicate demuestran que, en varias tareas, el mecanismo entre datapoints realmente se utiliza.
  \item La quadratic datapoint attention es la limitación central; a futuro se requieren retrieval, sparsity o representative points.
\end{enumerate}
\end{importantbox}

\subsection{Un marco de decisión para la ingeniería}

Para decidir si conviene adoptar un mecanismo al estilo NPT en una tarea nueva, pueden plantearse cinco preguntas en orden: primero, si los otros examples realmente aportan información adicional durante la inference; segundo, si es posible acceder a esos examples de forma legítima, segura y con baja latencia; tercero, si el lote o el retrieval set cubre los neighbors relevantes; cuarto, si un baseline común de tree, k-NN o retrieval ya es suficiente; quinto, si las pruebas de corruption e intervention pueden demostrar que el modelo no depende únicamente de sus propias features o de un shortcut.

\begin{knowledgebox}{Lista de verificación práctica}
Primero, implementar un baseline per-row sólido; después, incorporar explicit retrieval; luego, comparar k-NN fijo, learned retriever y full datapoint attention; por último, realizar row corruption, neighbor removal, label intervention, análisis de sensibilidad a la composición del lote y OOD evaluation. Solo cuando mejoren a la vez el desempeño y la evidencia sobre el mecanismo debe atribuirse la ganancia al razonamiento entre datapoints.
\end{knowledgebox}

\subsection{Lecturas adicionales}

Se recomienda leer en el orden «mecanismo--arquitectura--validación»: primero, volver a \emph{Attention Is All You Need} para comprender el causal training; después, leer la sección 2 y los apéndices del artículo de NPT de Kossen et al. para dominar los tensor shapes y el masking; a continuación, contrastar Deep Sets, Set Transformer, Neural Relational Inference y Neural Processes, comparando permutation symmetry, message passing, context conditioning y uncertainty assumptions. En la práctica de ingeniería conviene seguir de cerca la sparse attention y el retrieval-augmented modeling, porque atacan directamente el quadratic bottleneck de NPT sobre el conjunto de datos completo.

\teachervoice{El cierre del ponente no afirma que el problema esté resuelto; deja el scaling y las tareas más ricas como direcciones abiertas. La síntesis más prudente es esta: NPT demuestra que «entrenar un algoritmo de predicción que lee el conjunto de datos» funciona y aporta experimentos poco comunes a nivel de mecanismo; al mismo tiempo, convierte la cobertura de la recuperación, la complejidad computacional y la fuga de datos en problemas de sistema ineludibles.}

\end{document}
